\documentclass[11pt]{article}
\usepackage[margin=2.5cm]{geometry}
\usepackage{booktabs}

\title{FinLingo++ -- Final Metrics}
\author{24222437}
\date{\today}

\begin{document}
\maketitle

\section*{RQ1 -- Simplification quality (n=91)}
\begin{table}[h!]
\centering
\begin{tabular}{lrrl}
\toprule
Metric & Result & Target & Status \\
\midrule
SARI & 49.89 & $>40$ & Pass \\
BERTScore F1 & 0.895 & $>0.88$ & Pass \\
FK grade & 9.89 & $\leq 9$ & Fail \\
\bottomrule
\end{tabular}
\end{table}

\section*{RQ2 -- Risk classification (n=91)}
\begin{table}[h!]
\centering
\begin{tabular}{lrrl}
\toprule
Metric & Result & Target & Status \\
\midrule
Macro-F1 & 0.563 & $>0.75$ & Fail \\
Accuracy & 0.560 & -- & -- \\
\bottomrule
\end{tabular}
\end{table}

\section*{RQ3 -- Faithfulness (n=182)}
\begin{table}[h!]
\centering
\begin{tabular}{lrrl}
\toprule
Metric & Result & Target & Status \\
\midrule
Primary (retrieved chunk) precision$^*$ & 1.00 & $>0.90$ & Fail$^*$ \\
Primary (retrieved chunk) recall & 0.006 & $>0.80$ & Fail \\
Secondary (source clause) precision & 0.906 & $>0.90$ & Pass (marginal) \\
Secondary (source clause) recall & 0.637 & $>0.80$ & Fail \\
\bottomrule
\end{tabular}
\end{table}
\noindent $^*$Precision of 1.00 comes from a single positive prediction out of 182 records --
not meaningful, kept for completeness. See root-cause note below.

\section*{Retrieval}
\begin{table}[h!]
\centering
\begin{tabular}{lr}
\toprule
Metric & Result \\
\midrule
Recall@8 (pre-rerank) & 0.115 \\
Recall@5 (post-rerank) & 0.082 \\
\bottomrule
\end{tabular}
\end{table}

\section*{RQ4 -- Ablations and comparisons}
\begin{table}[h!]
\centering
\begin{tabular}{lrl}
\toprule
Comparison & Result & Note \\
\midrule
Reranker L-6 vs L-12 & L-12 wins MRR/nDCG/hit@1 & tied recall@5, $\sim$1.8x latency \\
Multi-verifier: DeBERTa v3 (deployed) & P 0.937 / R 0.648 & -- \\
Multi-verifier: ModernBERT & 0.0 / 0.0 & degenerate \\
FaithBench: v3 vs base & $+0.038$ balanced acc. & CI [0.006, 0.072], $p=0.02$ \\
FinBERT baseline (comparator) & Macro-F1 0.69 & vs deployed S5's 0.56 \\
RAGAS-style judge (comparator) & P 0.689 / R 0.560 & deployed verifier wins \\
SelfCheckGPT-style (comparator) & P 0.849 / R 0.495 & deployed verifier wins \\
\bottomrule
\end{tabular}
\end{table}
\noindent Full ablation suite (6 pipeline variants + 5 comparator baselines): complete.
Source: \texttt{reports/ablations.json}.

\section*{Adapter check}
\begin{table}[h!]
\centering
\begin{tabular}{ll}
\toprule
Check & Result \\
\midrule
Structural (LoRA modules present, correct target modules) & Pass (648 modules) \\
Behavioural probe (14 pairs) & Inconclusive (see note) \\
\bottomrule
\end{tabular}
\end{table}
\noindent The behavioural probe compares base vs fine-tuned outputs on 14 pairs, expecting a
score change above a fixed threshold. All 14 probes turned out to be cases the base model
already scores near 0 or 1 with high confidence, leaving no room for a fine-tuning effect to
show up -- so the test is inconclusive, not evidence against the adapter. The FaithBench
result above is the properly-powered version of the same check.

\section*{Notes}
\begin{itemize}
\item RQ3 primary result is low because the verifier only checks the top 5 retrieved
  chunks, and the correct chunk is in that top 5 only about 7\% of the time on this
  benchmark -- confirmed as a retrieval/benchmark issue, not a code defect, by testing the
  same models on a different benchmark (FinanceBench), where recall@5 is around 70\%.
\item RQ2 target (macro-F1 $>0.75$) sits alongside a requirement that this stage stay
  prompted-only, with no fine-tuning. This is a known tension in the current design.
\item All numbers in this document are taken directly from \texttt{reports/final\_results.json}
  and the underlying run files it references (\texttt{ablations.json},
  \texttt{multiverifier.json}, \texttt{reranker\_comparison.json},
  \texttt{verifier\_faithbench\_v3\_reproduction.json}, \texttt{adapter\_attachment\_proof.json}).
\end{itemize}

\end{document}
